\section{讲座定位：这不是“把 Transformer 变浅”那么简单}

标题里的“近浅层”（near-shallow）容易让人误以为课程只讨论减少 Transformer 层数。讲者实际提出的是一条更具体、也更受条件约束的路线。第一步，把标准 self-attention 中“先投影、再比较”的一部分工作改写成“直接变换向量比较”；第二步，用 generalized co-occurrence（广义共现）为 softmax layer 构造非随机 warm start；第三步，在 embedding 固定时缓存底层向量比较，使一部分计算不再随每次训练或推理重复；第四步，观察这种结构能否支持只用本地应用数据训练的 tiny controller。这里的 near-shallow 指\textbf{可复用的底层计算让有效在线路径更浅}，不是网络只剩一层。

\subsection{先修复来源：旧讲义的标题和内容都不可靠}

旧稿把课程日期写成 2026 年，使用不存在的 Stanford lecture URL，并加入 drift gate、attention observability board、fairness dashboard、weekly newsletter、RLHF governance roadmap 等大量内容。这些机制既不在官方课程页和录像中，也不在讲者相关论文里。本讲现以 2024 年 5 月 2 日课堂、Stanford Online 官方录像 \texttt{zL9B3eXq0gY}、官方 \texttt{en-US} 人工字幕和 27 个录像恢复教学 slide states 为唯一课堂主干。

\lecturefigure{slide-01-motivation.jpg}{核心问题：能否直接学习“比较结果到 feature weights”的变换}{Stanford Online 官方录像 00:04:39}

第一张教学 slide 对比两条路径。标准 attention 通过 $W_q,W_k$ 改变表示空间，再用 dot product 产生 score；讲者的替代形式保留原向量比较 $XX^\top$，另用矩阵 $W$ 把这些比较转换成 attention logits。后文所有显式初始化、缓存和 edge 实验都围绕这次改写展开。

\begin{importantbox}{本讲的三个核心问题}
\begin{enumerate}
  \item \textbf{Initialization：}能否不用随机参数起步，而从输入--输出统计量直接得到一个有用的初值？
  \item \textbf{Reusable computation：}哪些表示或向量比较可以冻结并缓存，从而少做重复计算？
  \item \textbf{Application-only learning：}一个小模型能否不依赖通用预训练，只从某个本地任务的交互数据学习？
\end{enumerate}
\end{importantbox}

\teachervoice{00:01:19--00:02:49，Williams 先说明自己的数学、物理与 quantitative linguistics 背景。他刻意从统计结构和显式公式切入，而不是从常见的 scaling recipe 切入。}

\subsection{证据分层：标准知识、论文推导与课堂实验不能混成一件事}

本讲同时出现标准 Transformer 公式、团队预印本中的推导、课堂 slide 上的实验结果、Le Potato 可行性演示，以及 Q\&A 中的未来设想。为避免把研究原型写成成熟配方，阅读时应保持下列分层。

\begin{knowledgebox}{证据账本}
\begin{tabularx}{\textwidth}{L{0.17\textwidth} L{0.29\textwidth} X}
\toprule
层次 & 例子 & 正确读法 \\
\midrule
标准背景 & softmax、dot-product attention、困惑度、padding & 可作为公共技术基础 \\
论文推导 & generalized co-occurrence 显式解、SAFFU hidden targets & 需遵守论文假设与近似条件 \\
课堂实验 & warm/cold 曲线、MNIST、BabyLM、switch task & 只支持对应设置下的观测 \\
讲者判断 & packing 不是正确 context model、PLM 命名 & 作为研究观点，不当作共识 \\
未来工作 & convolution warm start、multimodal、RLHF、公开实现 & 尚未完成，不能写成已验证能力 \\
\bottomrule
\end{tabularx}
\end{knowledgebox}

\begin{warningbox}{SAFFU 不是现成的标准架构}
Self-Attentive Feed-Forward Unit（SAFFU，自注意力前馈单元）和 Precision Language Model（PLM，精确语言模型）是讲者团队使用的研究术语。它们不是 PyTorch/Hugging Face 中已经标准化的模块，也不能仅凭这堂课推出“训练 LLM 不再需要 backpropagation”。
\end{warningbox}

\subsection{本章小结}

真实课程围绕非随机初始化、缓存与本地小模型展开，不是旧稿中的治理流程。本讲义会把每一条结果放回其具体假设：softmax、非负特征、固定 embedding、有限数据与早期 edge prototype。

